\documentclass[
%fontsize=10pt, % Base font size
%twoside=false, % If true, use different layouts for even and odd pages (in particular, if twoside=true, the margin column will be always on the outside)
%secnumdepth=2, % How deep to number headings. Defaults to 2 (subsections)
%abstract=true, % Uncomment to print the title of the abstract
]{kaohandt}

% Choose the language
\usepackage[spanish]{babel} % Load characters and hyphenation
\usepackage[spanish=spanish]{csquotes}	% Spanish quotes

\usepackage{tikz}

% Load packages for testing
\usepackage{blindtext}
%\usepackage{showframe} % Uncomment to show boxes around the text area, margin, header and footer
%\usepackage{showlabels} % Uncomment to output the content of \label commands to the document where they are used

\graphicspath{{images/}{./}} % Paths where images are looked for

% Load mathematical packages for theorems and related environments.
\usepackage{kaotheorems}

% Load the bibliography package
\usepackage{kaobiblio}
\addbibresource{report-template.bib} % Bibliography file

% Load the package for hyperreferences
\usepackage{kaorefs}


% Some shortcuts

% Real numbers, Natural numbers, etc.
\renewcommand{\P}{\mathbb{P}}
\newcommand{\E}{\mathbb{E}}
\newcommand{\R}{\mathbb{R}}
\newcommand{\Q}{\mathbb{Q}}
\newcommand{\Z}{\mathbb{Z}}
\newcommand{\C}{\mathbb{C}}

% (Conditional) expected value
\newcommand{\mean}[2][]{%
	\ifthenelse{\isempty{#1}}%
	{\mathbb{E}\!\left[#2\right]}%
	{\mathbb{E}\!\left[#2\,\middle|\,\mathcal{F}_{#1}\right]}%
}


%----------------------------------------------------------------------------------------

\begin{document}
	
	%----------------------------------------------------------------------------------------
	%	REPORT INFORMATION
	%----------------------------------------------------------------------------------------
	
	\title[]{Fórmula de Jensen}
	
	\author[LM]{Lorenzo Martinelli\thanks{Universidad de Buenos Aires} }
	
	\date{\today}
	
	%----------------------------------------------------------------------------------------
	%	TITLE AND ABSTRACT
	%----------------------------------------------------------------------------------------
	
	\maketitle
	
	\margintoc
	
	\begin{abstract}
		\noindent
		La idea de esta breve nota es, por un lado enunciar y probar la fórmula de Jensen y por otro, mostrar algunas aplicaciones al crecimiento de los ceros de una función entera.
	\end{abstract}
	
	%{\noindent\textbf{Keywords:} \LaTeX, Kao, handout, article, report}
	
	\medskip
	
	%----------------------------------------------------------------------------------------
	%	MAIN BODY
	%----------------------------------------------------------------------------------------
	
	\section{Introducción}
	
	El objeto de estudio en esta nota serán las funciones holormorfas en todos el plano también llamadas enteras. En particular, la pregunta que nos guiará es la siguiente.
	
	¿Cómo crecen en el infinito estas funciones? Acá aparece un hecho fundamental, mientras más \emph{grande} es una función mayor puede ser su cantidad de ceros. Esto de alguna manera ya los sabemos para los polinomios donde la cantidad de ceros de un polinomio $P$ de grado $d$ es precisamente $d$ que es además, el orden (polinomial) de crecimiento de este.\sidenote[][*-6]{Consecuencia del teorema fundamental del álgebra.}

	\section{Fórmula de Jensen}
	
	Antes de enunciarla formalmente, vamos a recordar el teorema de Cauchy (¡¿Cuál de todos?!) y dos corolarios que van a ser de suma utilidad. 
	
	\begin{theorem}[Fórmula integral de Cauchy]
		Sea $f : \Omega \to \C$ con $\Omega$ abierto y que contiene la clausura de un disco $D_R(z_0)$. Luego, 
		$$f(z) = \frac{1}{2\pi i} \int_{\partial D_R(z_0)} \frac{f(\xi)}{\xi - z}d\xi$$
	\end{theorem}
	
	\begin{proof}
		Fijado $z \in D_R(z_0)$ consideremos $\Gamma_{\varepsilon, \delta}$ \sidenote[][*-6]{ Consideramos $C$ al borde del disco de la hipótesis
		\begin{center}
			\includegraphics[width=6cm]{/imagenes/curva.png}\\
		\end{center}
		}. Donde $\varepsilon$ es el radio del agujero centrado en $z$ y $\delta$ el ancho del pasillo que escapa. Ahora, tenemos que
		$$\int_{\Gamma} \frac{f(\xi)}{\xi - z}d\xi = 0$$
		pues en el interior de $\Gamma$, notando $F(\xi) = \frac{f(\xi)}{\xi - z}$ es holomorfa. Por otro lado, como $F$ es continua tenemos que cuando $\delta$ tiende a 0 las integrales de ambos lados del pasillo se cancelan, por lo que nos queda \sidenote[][*-2]{El menos viene de orientar en sentido anithorario el agujero.}		
		$$ \int_{C} F(\xi)d\xi - \int_{C_\varepsilon} F(\xi)d\xi = 0.$$
		Veamos ahora que ocurre con el segundo termino,
		$$\int_{C_\varepsilon} F(\xi)d\xi = \int_{C_\varepsilon} \frac{f(\xi) - f(z)}{\xi - z} + \frac{f(z)}{\xi - z}d\xi,$$
		cómo $f$ es holomorfa el termino de la izquierda esta acotado y, cuando tomamos límite con $\varepsilon$ tendiendo a $0$ tenemos que este se va a 0. Por otro lado, podemos calcular explicitamente el segundo termino,
		$$ \int_{C_\varepsilon} \frac{f(z)}{\xi - z}d\xi =  \int_{0}^{2\pi} \frac{f(z) i \varepsilon e^{i \theta}}{\varepsilon e^{i\theta}}d\theta = 2\pi if(z)$$
		volviendo a lo que nos habia quedado antes, tenemos el resultado.  
		
	\end{proof}
	
	\begin{corollary}[Valor medio]
		Dada $f : D_R(z_0) \to \C$ holomorfa entonces
		$$f(z_0) = \frac{1}{2\pi} \int_{0}^{2\pi} f(z_0 + re^{i\theta})d\theta, \quad \text{con } \, 0 < r < R.$$ 
	\end{corollary}
	
	\begin{corollary}
		Dada $f : D_R(z_0) \to \C$ holomorfa, si $u = Re(f)$ entonces
		$$u(z_0) = \frac{1}{2\pi} \int_{0}^{2\pi} u(z_0 + re^{i\theta})d\theta, \quad \text{con } \, 0 < r < R.$$ 
	\end{corollary}
	
	\begin{proof}
		Basta tomar parte real al resultado del primer corolario.
	\end{proof}
	
	Con esto en mente, enunciemos la fórmula de Jensen y probemosla.
	
	\marginnote{
	\begin{center}
		\includegraphics[width=4cm]{/imagenes/jensen.jpg}
		\captionof{figure}{Johan Jensen (8 de mayo de 1859).}
	\end{center}}
	
	\begin{theorem}[Fórmula de Jensen]
		Sea $f : \Omega \to \C$ holomorfa con $\Omega$ abierto, que contiene la clausura de un disco $D_R(0)$, $f(0) \neq 0$ y $f$ no se anula en el borde del disco. Si $z_1, ..., z_N$ son los ceros de $f$ en el interior del disco entonces
		$$\log|f(0)| = \sum_{i=1}^{N}{\log{\frac{|z_k|}{R}}} + \frac{1}{2\pi}\int_{0}^{2\pi} \log|f(Re^{i\theta})|d\theta.$$
	\end{theorem}
	
	\begin{proof}
		La prueba se divide en varios pasos.
		\begin{enumerate}
			\item Comencemos notando que si $f_1, f_2$ satisfacen las hipótesis del teorema entonces $f_1f_2$ también las satisface. Más aún satisface la fórmula pues el logaritmo separa sumas en productos y la unión de los ceros de $f_1, f_2$ son los ceros de $f_1f_2.$
			
			\item Sabemos que podemos escribir $f(z) = (z-z_1)\dots(z-z_N)g(z)$ con $g$ holomorfa y no nula en la clausura de $D_R$. Luego, basta ver el resultado para funciones como $g$ y de la forma $z - w$ con $w \in D_R.$
			
			\item Sea $g$ holomorfa que no se anula, queremos ver que
			$$\log{|g(0)|} = \frac{1}{2\pi} \int_{0}^{2\pi} \log{|g(Re^{i\theta})|}d\theta $$
			cómo $g$ no se anula, podemos definir en un disco un póco más grande $h$ holomorfa de manera que $e^{h(z)} = g(z).$\sidenote{Queda pendiente probarlo para cuando terminemos la demostración, pero esencialmente se deduce de poder definir una rama del logaritmo.} Ahora, $|g(z)| = |e^{h(z)}| = e^{Re(h(z))}.$ Luego, $\log{|g(z)|} = Re(h(z))$ y usando el corolario que vimos antes tenemos el resultado. 
			
			\item Nos queda ver el resultado para $f(z) = z - w$ con $w \in D_R.$ Tenemos que ver que
			$$\log{|w|} = \log{|f(0)|} = \log{\frac{|w|}{R}} + \frac{1}{2\pi} \int_{0}^{2\pi} \log{|Re^{i\theta} - w|}d\theta.$$ 
			Pero esto es,
			$$\log{|w|} = \log{|w|} - \log{R} + \frac{1}{2\pi} \int_{0}^{2\pi} \log{|Re^{i\theta} - w|}d\theta,$$
			mirando el segundo termino,
			$$ \frac{1}{2\pi} \int_{0}^{2\pi} \log{|Re^{i\theta} - w|}d\theta = \log{|R|} + \int_{0}^{2\pi} \log{|e^{i\theta} - \frac{w}{R}|}d\theta.$$
			Luego, basta probar que
			$$\int_{0}^{2\pi} \log{|e^{i\theta} - a|}d\theta = 0, \quad \text{con } |a| < 1$$
			Que es equivalente, via cambio de variables $\theta = -\theta$ a
			$$\int_{0}^{2\pi} \log{|1 - ae^{i\theta}|}d\theta = 0.$$
			De la misma manera que antes, llamando $F(z) = 1 - az$, existe $G(z)$ tal que $F(z) = e^{G(z)}$ y, entonces $\log|F(0)| = Re(G(z))$, luego por la fórmula de valor medio tenemos el resultado.
		\end{enumerate}
		
	\end{proof}
	
	Antes de continuar vamos a enunciar y probar un pequeñp resultado que usamos en la demostración,
	
	\begin{proposition}
		Sea $f: \Omega \to \C$ holomorfa, con $\Omega$ simplemente conexo. Si $f$ nunca se anula en su domino entonces existe $g(z)$ de manera tal que $e^{g(z)} = f(z).$ 
	\end{proposition}
	
	\begin{remark}
		Notemos que antes $\Omega$ no era simplemente conexo pero el disco sí.
	\end{remark}
	
	\begin{proof}
		Fijado $z_0 \in \Omega$ consideremos $c_0$ de manera tal que $e^{c_0} = f(z_0)$. Así, podemos definir
		$$ g(z) = \int_{\gamma} \frac{f'(z)}{f(z)}dz + c_0$$
		con $\gamma$ curva de $z_0$ a $z$. Notemos que está bien definida pues $\Omega$ es simplemente conexo y $f$ es holomorfa. Ahora, derivando $f(z)e^{-g(z)}$ tenemos que es $0$ y luego,
		$f(z) = Ke^{g(z)}.$ Evaluando en $z_0$ tenemos que $K = 1$.
	\end{proof}
	
	\subsection{Aplicaciones}
	
	Si notamos $n(r)$ a la cantidad de ceros en el disco de centro $0$ y radio $r$ podemos ver, con ayuda de un lema, que
	$$\int_{0}^{R}n(r)\frac{dr}{r} = \frac{1}{2\pi} \int_{0}^{2\pi} \log{|f(Re^{i\theta})|}d\theta - \log{|f(0)|},$$
	es decir, la relación entre los ceros en el disco y el crecimiento del radio esta controlado por el promedio logarítmico de $f$.
	
	\begin{lemma}
		Si $z_1, \dots,z_N$ son los ceros de $f$ dentro de $D_R$ entonces
		$$\int_{0}^{R}n(r)\frac{dr}{r} = \sum_{k=1}^{N} \log{\frac{R}{|z_k|}}.$$
	\end{lemma}
	
	\begin{proof}
		Primero notemos que
		$$\sum_{k=1}^{N} \log{\frac{R}{|z_k|}} = \sum_{k=1}^{N} \int_{|z_k|}^{R}n(r)\frac{dr}{r}.$$
		Si definimos ahora $n_k(r)$ la indicador del conjunto $[|z_k|, +\infty)$ podemos ver que $\sum_{k=1}^{N} n_k(r) = n(r).$ En consecuencia, podemos ver que
		$$\sum_{k=1}^{N}\int_{|z_k|}^{R} \frac{dr}{r} = \sum_{k=1}^{N}\int_{0}^{R} n_k(r)\frac{dr}{r} = \int_{0}^{R} \sum_{k=1}^{N} n_k(r)\frac{dr}{r} = \int_{0}^{R} n(r)\frac{dr}{r}.$$
	\end{proof}
	
	\subsubsection{Funciones de orden finito}
		
	Sea $f$ una función entera. Si existe $\rho > 0$ y constantes $A, B > 0$ tales que 
	$$ |f(z)| \leq Ae^{B|z|^\rho}$$
	decimos que $f$ tiene orden menor o igual a $\rho.$ Por otro lado, definimos el orden de $f$ como 
	$$\rho_f = \inf{\rho}.$$
	
	\begin{theorem}
		Si $f$ es entera con orden menor o igual a $\rho$. Luego,
		\begin{enumerate}
			\item Existen $C > 0$ y $r$ suficientemente grande tal que $n(r) \leq Cr^{\rho}.$
			\item Si $z_1, z_2 \dots$ son los ceros de $f$, con $z_k \neq 0$ entonces para $s > \rho$ vale que 
			$$ \sum_{k\geq1} \frac{1}{|z_k|^{s}} < \infty.$$
		\end{enumerate}
	\end{theorem}  
	
	\begin{proof}
		
		Probemos entonces ambos items en orden.
		
		\begin{enumerate}
			\item Comencemos notando que alcanza probar el resultado para $f(0) \neq 0$ pues podemos considerar $\frac{f(z)}{z^{l}}$ si $l$ es la multiplicidad de $0$ como raíz, los ceros de esta función difieren solo en $l$ de los de $f$. Si $f(0) \neq 0,$ podemos usar el lema de antes y tenemos
			$$\int_{0}^{R} n(x)\frac{dx}{x} = \frac{1}{2\pi} \int_{0}^{2\pi} \log{|f(Re^{i\theta})|}d\theta - \log{|f(0)|}$$
			Eligiendo $R = 2r$,
			$$\int_{r}^{2r} n(x)\frac{dx}{x} \leq \frac{1}{2\pi} \int_{0}^{2\pi} \log{|f(2re^{i\theta})|}d\theta - \log{|f(0)|}.$$
			Además como $n$ es creciente, 
			$$ \log{2}n(r) = \int_{r}^{2r} n(r) \frac{dx}{x} \leq \int_{r}^{2r} n(x)\frac{dx}{x}.$$
			Solo nos queda usar que $f$ tiene orden menor o igual a $\rho$,
			$$\frac{1}{2\pi} \int_{0}^{2\pi} \log{|f(2re^{i\theta})|}d\theta \leq \frac{1}{2\pi} \int_{0}^{2\pi} \log{Ae^{i(2r)^{\rho}}}d\theta = C'r^{\rho}.$$
			
			\item Notemos que solo pueden existir finitas raices con módulo menor a 1 pues sino, por identidad, f sería nula. Luego,
			\[
			\begin{aligned}
			\sum_{|z_k|\geq1} \frac{1}{|z_k|^{s}} &= \sum_{j\geq 0} \sum_{2^{j} \leq |z_k| \leq 2^{j+1}} |z_k|^{-s} \\
			&\leq \sum_{j\geq 0} n(2^{j+1})2^{-js} \\
			&\leq \sum_{j\geq 0} C2^{(j+1)\rho}2^{-js} \\
			&= \sum_{j\geq 0} C'2^{j(\rho - s)} < \infty
			\end{aligned} 
			\]
			Donde el último termino es finito pues $s > \rho.$
		\end{enumerate}
	\end{proof}
	
	Para terminar, veamos un ejemplo que muestra que la condición de $s > \rho$ no puede ser mejorada.
	
	\begin{example}
		Consideremos $f(z) = \sin{z\pi}$, tenemos en este caso que sus raices son de la forma $z_k = k$. Además, podemos ver que $|f(z)| \leq e^{|z|\pi}$, es decir, que tiene orden menor igual a 1. Precisamente, $\sum_{k\geq1} \frac{1}{k^s}$ es finito si y sólo si $s > 1.$ 
	\end{example}
\end{document}